\begin{tikzpicture}[font=\small,>=stealth]
\draw[rounded corners=3pt,gray!45] (0,0) rectangle (14.2,17);
\node at (7.1,16.45) {不用准线：分别看两个焦点三角形，夹角的余弦变号};
\node at (7.1,15.35) {上图只看 $F_1FA$：另一条边由椭圆定义补成 $2a-r$};\begin{scope}[shift={(4.4,11.75)},scale=.65]
\draw[teal!60!black] (0,0) ellipse (5 and 3);
\fill[blue!75!black,opacity=.07] (-4,0)--(4,0)--(4.6428571429,1.1134612334)--cycle;\draw[gray,thick] (-4,0)--(4,0);\draw[blue!75!black,very thick] (4,0)--(4.6428571429,1.1134612334);\draw[violet,thick] (-4,0)--(4.6428571429,1.1134612334);
\fill (-4,0) circle(2pt) node[below=8pt] {$F_1$};\fill (4,0) circle(2pt) node[right=22pt] {$F$};\fill (4.6428571429,1.1134612334) circle(2pt) node[above=9pt] {$A$};
\draw[violet] (4.4,.692820323) arc[start angle=60,end angle=180,radius=.8];\node[violet] at (2.5,-1.3) {$\pi-\theta$};\draw[gray,thin] (2.9,-.95)--(3.5,.35);
\node[below=10pt] at (0,0) {$2c$};\end{scope}
\node[anchor=west,align=left] at (9.2,11.75) {$AF=r$\\[12pt]$AF_1=2a-r$\\[15pt]夹角 $\pi-\theta$\\[12pt]余弦 $-\cos\theta$};
\begin{scope}[shift={(4.4,5.55)},scale=.65]
\draw[teal!60!black] (0,0) ellipse (5 and 3);
\fill[orange!85!black,opacity=.07] (-4,0)--(4,0)--(2.5000000000,-2.5980762114)--cycle;\draw[gray,thick] (-4,0)--(4,0);\draw[orange!85!black,very thick] (4,0)--(2.5000000000,-2.5980762114);\draw[violet,thick] (-4,0)--(2.5000000000,-2.5980762114);
\fill (-4,0) circle(2pt) node[below=8pt] {$F_1$};\fill (4,0) circle(2pt) node[right=22pt] {$F$};\fill (2.5000000000,-2.5980762114) circle(2pt) node[below=9pt] {$B$};
\draw[violet] (3.2,0) arc[start angle=180,end angle=240,radius=.8];\node[violet] at (2.8,-.45) {$\theta$};
\node[above=8pt] at (0,0) {$2c$};\end{scope}
\node[anchor=west,align=left] at (9.2,5.55) {$BF=s$\\[12pt]$BF_1=2a-s$\\[15pt]夹角 $\theta$\\[12pt]余弦 $+\cos\theta$};

\node at (7.1,8.7) {下图只看 $F_1FB$：另一条边补成 $2a-s$};
\node at (7.1,1.6) {两图同一椭圆、同一条弦、同一尺度；只切换正在使用的三角形。};
\node at (7.1,.6) {图取 $\theta=60^\circ$；一般 $0<\theta<\pi$ 时，两个夹角仍分别为 $\pi-\theta,\theta$。};
\end{tikzpicture}
